% !TEX root = ../main.tex

\section{Theoretical Background (total: 5-6.5 pages)}
Impulse response functions can be nonparametically defined as the difference between two conditional forecasts of the outcome variable (e.g. Koop (1996)). Both forecasts are conditioned on the same history up to period $t-1$, but differ in the realization at time $t$. While one conditions on a shock of size $\delta$ to the variable of interest, the other one does not. Following \textcite{ramey2016macroeconomic} shocks hereby refers to a significant exogenous force representing unanticipated movements in exogeneous variables, that is uncorrelated with other current and lagged control variables and other exogeneous shocks. The observed difference in the projected paths at horizon $h$ then identifies the causal effect of the shock in population:

\begin{equation}
    \mathrm{IR}(h,\delta) \equiv \mathbb{E}\left[y_{t+h}|s_t = s_0 + \delta; \mathbf{x_t}\right] - \mathbb{E}\left[y_{t+h}|s_t = s_0; \mathbf{x_t}\right]
\end{equation}

Vector Autoregressions and Local Projections serve as two complementary  approaches to estimating impulse responses in finite samples. The former,  introduced by \textcite{sims1980macroeconomics}, parametrizes the joint  dynamics of all variables in a single autoregressive system and impulse  responses are then obtained as iterated forecasts from the estimated  companion matrix. The latter, established by \textcite{jorda2005estimation},  takes a direct forecasting approach by estimating the responses at each horizon independently via OLS, projecting the future outcome directly onto current and lagged values of the data. The two approaches compound estimation error differently across horizons and, at equal lag order, occupy opposite ends of a bias-variance spectrum in finite samples \parencites{li2024local}{olea2025local}. Local projections tend toward  higher variance, vector autoregressions toward higher bias under misspecification. 
% For unrestricted lag lengths, \textcite{plagborg2021local} establish that both methods recover the exact same impulse response function in population. 
More recently, however, \textcite{ludwig2026local} shows that the equal-order comparison is itself misleading. At a given lag order $p$, LP$(p)$ is strictly more general than VAR$(p)$, so the conventional bias-variance trade-off conflates the choice of estimator with the choice of model complexity.


\subsection{Introduction to VARS (1.5 - 2 pages)}

Following \textcite{lutkepohl2005}, we start by defining a zero mean VAR(p) population model of the form where the outcome of interest, $y_t$, is a $K \times 1$ random vector, $A_i$ are fixed $K\times K$ coefficient matrices and $u_t$ are a K dimensional white noise processes with $\mathbb{E}\left[ u_t\right] = 0$, $\mathbb{E}\left[ u_tu_t{'}\right] = \Sigma_u$, i.e. being a nonsingular covariance matrix, and $\mathbb{E}\left[ u_tu_s{'}\right] = 0$ for $t\neq s$. Instead of writing out the full lag structure model the model, we can rewrite the full VAR(p) in its corresponding VAR(1) companion form such that
%
\begin{equation}
    Y_t = \mathbf{A} Y_{t-1} + U_t
\end{equation}
%
where
%
\begin{equation*}
Y_t := \begin{bmatrix} y_t \\ y_{t-1} \\ \vdots \\ y_{t-p+1} \end{bmatrix} _{Kp\times 1}, 
\mathbf{A} := \begin{bmatrix} 
A_1 & A_2 & \cdots & A_{p-1} & A_p \\ 
I_K & 0   & \cdots & 0       & 0   \\
0   & I_K &        & 0       & 0   \\
\vdots & & \ddots & & \vdots \\
0   & 0   & \cdots & I_K     & 0
\end{bmatrix} _{Kp\times Kp},
U_t := \begin{bmatrix} u_t \\ 0 \\ \vdots \\ 0 \end{bmatrix} _{Kp\times 1}.
\end{equation*}
%
Assuming $Y_t$ satisfies the stability condition such that all roots, $z$, of the characteristic equation lie strictly outside the complex unit circle, the process is covariance stationary and admits a Wold moving-average representation:
%
\begin{equation}
det(I_{Kp} - \mathbf{A}z) \neq 0 \text{ for } |z| \le 1 \iff Y_t = \sum_{l=0}^{\infty}\mathbf{A}^{l}U_{t-l}.
\end{equation}
%
Restricting attention to the top $K$ rows of the Wold MA representation yields the moving-average representation for $y_t$:
%
\begin{equation}
    y_t = \sum_{\ell=0}^{\infty} \Psi_\ell\, u_{t-\ell}, 
\qquad \Psi_\ell := [\mathbf{A}^\ell]_{1:K,\,1:K}
    \label{eq:wold_yt}
\end{equation}
%
where $\Psi_h$ is the reduced-form impulse response function at horizon $h$, measuring the response of $y_t$ to a unit reduced-form innovation $u_{t-h}$. It corresponds to the upper-left $K \times K$ block of $\mathbf{A}^h$, which follows directly from the fact that $U_{t-i} = (u_{t-i}, 0, \ldots, 0)'$ has nonzero entries only in its first $K$ positions. Since the components of $u_t$ are generally contemporaneously correlated, $\Sigma_u = \mathbb{E}(u_t u_t') \neq I_K$, a unit shock to one component of $u_t$ will in general be accompanied by simultaneous movements in the other components. Consequently, the reduced-form impulse responses $\Psi_h$ do not admit a structural interpretation. Structural identification is achieved by imposing restrictions on  the impact matrix $B$ in $u_t = B\varepsilon_t$, where $\varepsilon_t \sim (0, I_K)$ are mutually uncorrelated structural shocks, so that $\Sigma_u = BB'$.


The most common approach is recursive Cholesky identification \parencite{ramey2016macroeconomic}, which restricts $B$ to be lower triangular with positive diagonal entries, uniquely recovering $B$ as the Cholesky factor of $\Sigma_u$. The recursive ordering imposes a causal structure: the variable ordered first responds contemporaneously to all shocks, while the variable ordered last is insulated from all  contemporaneous feedback. The diagonal entries $b_{ii}$, $i = 1, \ldots, K$, measure the direct impact of structural shock $i$ on variable $i$, while the off-diagonal entries $b_{ij}$, $1 \leq j < i \leq K$, capture the contemporaneous spillover from structural shock $j$ to variable $i$.
Substituting $u_t = B\varepsilon_t$ into \eqref{eq:wold_yt} yields the structural moving-average representation
%
\begin{equation}
    y_t = \sum_{\ell=0}^{\infty} \Theta_\ell\, \varepsilon_{t-\ell}, 
\qquad \Theta_\ell := \Psi_\ell B, 
    \qquad B = \begin{pmatrix} 
b_{11} & 0      & \cdots & 0      \\ 
b_{21} & b_{22} & \cdots & 0      \\ 
\vdots &        & \ddots & \vdots \\ 
b_{K1} & b_{K2} & \cdots & b_{KK} 
\end{pmatrix}
    \label{eq:sma}
\end{equation}
%
where the $(i,j)$-th element of $\Theta_h$ measures the response of 
variable $i$ to a unit innovation in structural shock $j$ at horizon $h$. 
The structural impulse response of variable $i$ to shock $j$ is thus
%
\begin{equation}
    \theta_h^{ij} = [\Psi_h B]_{ij}
    \label{eq:sirf}
\end{equation}
%
where $[\,\cdot\,]_{ij}$ denotes the $(i,j)$-th element of a matrix.

Estimation of the structural impulse response $\hat{\theta}_h^{ij}$ 
proceeds in two stages. In the first stage, the VAR coefficient matrices 
$\hat{A} = (\hat{A}_1, \ldots, \hat{A}_p)$ are estimated equation-by-equation 
via OLS, which is efficient under the stability condition and yields
%
\begin{equation}
    \sqrt{T}\,\bigl(\mathrm{vec}(\hat{A}) - \mathrm{vec}(A)\bigr) 
    \xrightarrow{d} \mathcal{N}\!\left(0,\; \Gamma^{-1} \otimes \Sigma_u\right),
\end{equation}
%
where $\Gamma = \mathbb{E}(Z_t Z_t')$ is the second-moment matrix of the 
stacked regressors $Z_t = (1, y_{t-1}', \ldots, y_{t-p}')'$, and the 
innovation covariance $\Sigma_u$ is replaced by its residual analogue 
$\hat{\Sigma}_u$ \parencite{lutkepohl2005}. %Lag order selection follows standard information criteria \textcite{akaike1974new,schwarz1978estimating}.

In the second stage, structural objects are recovered from the 
reduced-form estimates. The impact matrix $\hat{B}$ is the Cholesky 
factor of $\hat{\Sigma}_u$, and reduced-form impulse responses are 
obtained via the recursion $\hat{\Psi}_h = \sum_{\ell=1}^{\min(h,p)} 
\hat{A}_\ell\,\hat{\Psi}_{h-\ell}$, $\hat{\Psi}_0 = I_K$. The 
structural impulse response estimator is then a smooth nonlinear function of the first-stage estimates.
%
\begin{equation}
    \hat{\theta}_h^{ij} = e_i'\,\hat{\Psi}_h\,\hat{B}\,e_j.
\end{equation}
%

Since $\hat{\theta}_h^{ij} = g(\hat{\alpha},\hat{\sigma})$ is 
differentiable, the delta method yields
%
\begin{equation}
    \sqrt{T}\,\bigl(\hat{\theta}_h^{ij} - \theta_h^{ij}\bigr)
    \xrightarrow{d}
    \mathcal{N}\!\left(0,\;
    \frac{\partial g}{\partial\alpha'}\,\Sigma_{\alpha}\,
    \frac{\partial g'}{\partial\alpha}
    +
    \frac{\partial g}{\partial\sigma'}\,\Sigma_{\sigma}\,
    \frac{\partial g'}{\partial\sigma}
    \right),
\end{equation}
%
where $\alpha = \mathrm{vec}(A)$, $\sigma = \mathrm{vech}(\Sigma_u)$, 
and the Jacobian $\partial\,\mathrm{vec}(\Psi_h)/\partial\alpha'$ 
admits a closed-form expression \parencite{lutkepohl1990asymptotic}. The resulting confidence intervals are valid pointwise in $h$ under stationarity, but not uniformly over the persistence of the process, motivating a direct examination of finite-sample behaviour.

The favourable asymptotic properties established above are qualified by three interrelated problems that arise in the sample sizes typical of macroeconomic applications.
that arise in the sample sizes typical of macroeconomic applications.
First, estimation error propagates through iteration. Because 
$\hat{\Psi}_h$ is the upper-left block of $\hat{\mathbf{A}}^h$, error in 
$\hat{A}$ enters the impulse response through powers of the companion 
matrix and is compounded multiplicatively as the horizon grows. Even small 
coefficient biases therefore translate into substantial distortions of the 
estimated responses at longer horizons \parencite{kilian1998small}. In the 
limiting case of unit or near-unit roots, \textcite{phillips1998impulse} 
shows that long-horizon impulse responses from unrestricted VARs are not 
even consistently estimated, converging instead to random variables rather 
than the true responses.
Second, the autoregressive coefficients themselves are biased downward in 
short samples. It is a long-established result that least squares 
underestimates the persistence of an autoregressive process 
\parencites{marriott1954bias}{nicholls1988bias}{pope1990biases}, and because $\hat{\Psi}_h$ inherits this bias through the recursion, the estimated impulse responses decay too quickly and understate the true persistence of the system. The distortion is most severe precisely when persistence is high and the sample is short. 
Third, these distortions undermine delta-method inference. When the largest roots approach the unit circle, the asymptotic normal approximation deteriorates and the resulting confidence intervals undercover, increasingly so at longer horizons \parencites{kilian1998small}{kilian1999finite}. 

Furthermore, a distinct class of distortions arises when the 
autoregressive structure is itself misspecified: if the true process has 
moving-average components, a finite-order VAR is generically 
inconsistent regardless of sample size, and while higher-order AR 
approximations can asymptotically absorb the omitted MA dynamics, the 
required lag expansion rapidly exhausts degrees of freedom in samples of 
typical macroeconomic length \parencite{braun1993misspecifications}. The 
consequences of lag-order choice are strongly asymmetric---overfitting 
merely inflates variance, whereas underfitting eliminates higher-order 
dynamics of the response curve and produces spuriously tight intervals 
with coverage far below nominal levels \parencite{kilian2001impulse}. 
Crucially, both channels are population-level phenomena: 
\textcite{ravenna2007vector} shows that truncation bias in the VAR 
coefficients and the resulting contamination of the structural rotation 
matrix persist even as $T \to \infty$.


\subsection{Introduction to Local Projections (1.5 -2 pages)}

Whereas the VAR recovers the impulse response by iterating a single fitted
model forward through powers of the companion matrix, the local projection
estimates each horizon's response directly from a separate regression
\parencite{jorda2005estimation}. Following \textcite{jorda2005estimation} and
\textcite{jorda2025local}, fix a horizon $h$ and project the future outcome
$y_{t+h}$ onto the current observation $y_t$, controlling for the $p$ lags
that span the conditioning information of the system,
%
\begin{equation}
    y_{t+h} = \nu_h + \Phi_h\, y_t
            + \sum_{\ell=1}^{p} \Xi_{h,\ell}\, y_{t-\ell}
            + \xi_{t+h}^{(h)},
    \qquad h = 0, 1, \ldots, H,
    \label{eq:lp_reduced}
\end{equation}
%
where $\nu_h$ is a $K \times 1$ vector of constants, $\Phi_h$ and
$\Xi_{h,\ell}$ are $K \times K$ coefficient matrices, and $\xi_{t+h}^{(h)}$
is the horizon-$h$ projection residual. The collection of $H+1$ separate
regressions in \eqref{eq:lp_reduced} is what \textcite{jorda2005estimation} terms the local projection. Each is estimated in isolation, so the lag length $p$ need not be common across horizons and no parametric model of the propagation
mechanism is imposed.
 
The object of interest is the coefficient on the contemporaneous regressor.
Partialling the $p$ lags out of $y_t$ leaves the one-step reduced-form
innovation $u_t = y_t - \mathbb{E}\!\left[y_t \mid y_{t-1}, \ldots,
y_{t-p}\right]$, so by the Frisch--Waugh--Lovell theorem $\Phi_h$ equals the population projection coefficient of $y_{t+h}$ on $u_t$. Inserting the Wold representation \eqref{eq:wold_yt}, $y_{t+h} = \sum_{\ell=0}^{\infty}
\Psi_\ell\, u_{t+h-\ell}$, and using that $u_t$ is orthogonal to every
$u_{t+h-\ell}$ with $\ell \neq h$, the projection isolates a single term,
%
\begin{equation}
    \Phi_h = \mathbb{E}\!\left[y_{t+h}\, u_t'\right]
             \bigl(\mathbb{E}\!\left[u_t u_t'\right]\bigr)^{-1}
           = \Psi_h,
    \qquad \Phi_0 = \Psi_0 = I_K.
    \label{eq:lp_psi}
\end{equation}
%
The local projection slope therefore recovers exactly the reduced-form
impulse response $\Psi_h$ of \eqref{eq:wold_yt}, but reads it off a direct
regression rather than from the recursion $\hat{\Psi}_h = \sum_{\ell=1}^{\min(h,p)}
\hat{A}_\ell \hat{\Psi}_{h-\ell}$.
 
Structural identification proceeds exactly as in the VAR case and is
logically distinct from the choice of estimator
\parencite{plagborg2021local}. Any identification scheme selects the
structural shock as a fixed linear combination of the reduced-form
innovations, $\varepsilon_{j,t} = \bigl(B^{-1}\bigr)_{j\cdot}\, u_t$
\parencite[p.~967]{plagborg2021local}, so the structural response is obtained
by post-multiplying the reduced-form responses by the impact matrix $B$,
%
\begin{equation}
    \theta_h^{ij} = [\Phi_h B]_{ij} = [\Psi_h B]_{ij}
                  = e_i'\, \Psi_h\, B\, e_j,
    \label{eq:lp_sirf}
\end{equation}
%
which is the same estimand as \eqref{eq:sirf}. Under recursive
(Cholesky) identification $B$ is the lower-triangular factor of $\Sigma_u$,
and \eqref{eq:lp_sirf} can equivalently be obtained in a single step by
ordering the variables and regressing $y_{i,t+h}$ on the contemporaneous
value of the shock-bearing variable, controlling for the variables ordered
before it together with the lags \parencites{jorda2005estimation}{plagborg2021local}.
 
The same construction accommodates the remaining standard schemes. Recursive, long-run, and sign restrictions are each implementable through the projection \eqref{eq:lp_reduced} because all of them merely fix the vector $b$ that maps the Wold innovations into the shock of interest
\parencite[p.~967]{plagborg2021local}. When an external instrument or proxy for the shock is available, ordering the instrument first in the conditioning set---the ``internal instrument'' approach---identifies the relative structural responses even when the shock is noninvertible, whereas the external SVAR-IV/LP-IV implementation of \textcite{stock2018identification} requires invertibility \parencite[p.~973]{plagborg2021local}. Because identification operates on the population innovations rather than on the estimator, any scheme available to the VAR can be carried out with local projections and vice versa \parencite[p.~975]{plagborg2021local}.
 
Each regression in \eqref{eq:lp_reduced} is estimated equation-by-equation by OLS, and the structural estimator $\hat{\theta}_h^{ij} = e_i'\, \hat{\Phi}_h\, \hat{B}\, e_j$ is, like its VAR counterpart, a smooth function of the first-stage coefficients and the residual covariance. Because the projection coefficient \emph{is} the impulse response, no companion-matrix recursion intervenes between estimation and the reported response
\parencite[p.~65]{jorda2025local}. Inference, however, is complicated by the overlapping-horizon structure of the residual: iterating \eqref{eq:lp_reduced} shows that $\xi_{t+h}^{(h)}$ is a moving average of the one-step forecast errors accruing between $t$ and $t+h$, and is therefore serially correlated of order $h-1$ even when the innovations are white noise \parencite[p.~163]{jorda2005estimation}. Accordingly \textcite{jorda2005estimation} recommends a
heteroskedasticity- and autocorrelation-consistent (HAC) covariance
estimator, such as Newey--West \parencite[p.~75]{jorda2025local}, which
delivers pointwise-valid, asymptotically normal inference under stationarity, but---mirroring the delta-method intervals of the VAR---neither uniformly over the persistence of the process nor at horizons that grow with the sample.
 
\textcite{montiel2021local} show that a minimal modification removes the need for the HAC correction. Augmenting the projection with one additional lag renders the effective regressor of interest stationary even under a unit root; in the leading scalar case the lag-augmented projection
%
\begin{equation}
    y_{t+h} = \beta_h\, y_t + \tilde{\beta}_h\, y_{t-1} + \xi_{t+h}^{(h)}
    \label{eq:lp_aug}
\end{equation}
%
has a coefficient on $y_t$ that equals the coefficient on the residualized
innovation $u_t$ in a regression on $(u_t, y_{t-1})$. The regression score
$\xi_{t+h}^{(h)} u_t$ is then serially uncorrelated---because $u_t$ is
uncorrelated with its own leads and lags---so that, contrary to conventional
wisdom, the ordinary heteroskedasticity-robust Eicker--Huber--White standard
error is sufficient despite the serially correlated residual
\parencite[pp.~1794--1795]{montiel2021local}. The resulting $t$-statistic is
asymptotically standard normal \emph{uniformly} over $\rho \in [-1,1]$ and
over horizons satisfying $h/T \to 0$, dispensing with the bandwidth and kernel
choices of HAC inference and restoring exactly the robustness to persistence
and long horizons that VAR inference lacks
\parencite[p.~1809]{montiel2021local}; \parencite[p.~76]{jorda2025local}.
 

The favourable and robust properties just described are, as for the VAR,
qualified in the sample sizes typical of macroeconomic applications, and the qualifications run largely opposite to those of the autoregression.
 
First, the price of estimating each horizon by an unrestricted regression is variance. The local projection discards the cross-equation and cross-horizon restrictions that the VAR exploits through its recursion, an efficiency loss that is increasing in the horizon \parencite{li2024local}. Since the two estimators target the same response in population \parencite[p.~975]{plagborg2021local},
no method dominates in mean squared error across data-generating processes; in
finite samples the comparison reduces to a bias--variance trade-off in which
the VAR is typically more precise at short horizons and the local projection
less biased at long ones \parencites{li2024local}{montiel2021local}.
 
Second, the local projection is not unbiased in finite samples.
\textcite{herbst2024bias} show that when the data are persistent---the case of
interest for impulse-response analysis---LP point estimates can be severely
biased at the sample lengths common in the literature, that the bias links the
horizons to one another, and that it admits a simple analytical approximation
and bias correction. The near-unit-root downward bias familiar from
autoregressive estimation thus re-emerges directly in the projection
coefficient \parencite[p.~75]{jorda2025local}.
 
Third, and offsetting these costs, the single-equation structure confers a
robustness the VAR lacks. Because the response is estimated directly rather
than by iterating a fitted model, specification errors do not compound across
horizons: \textcite{jorda2024long} show that in infinite-order
processes the local projection has lower bias than the VAR at horizons beyond the truncation lag, precisely because the coefficient is estimated without propagating coefficient error through powers of the companion matrix, which also makes the projection more forgiving of lag-length misspecification \parencite[p.~65]{jorda2025local}. This robustness is not unqualified on the inference side, however: \textcite{herbst2024bias} caution that in the same small, persistent samples the autocorrelation-robust standard errors of the classical approach can substantially understate sampling uncertainty---a further argument for the lag-augmented, heteroskedasticity-robust inference of \textcite{montiel2021local}.


\subsection{On the Equivalence of VARs and Local Projections (2-2.5 pages)}

As Local Projection inference gained in popularity over the last couple of years, more and more scholars have started comparing the asymptotical and finite sample trade-offs of both methods. Prior research has established the notion that SVARs are computational more efficient while LPs inherit favorable robustness to misspecification (see e.g.\cites{ramey2016macroeconomic}{jorda2005estimation}). The notion applies in analogy to the comparison of iterated vs. direct forcasting as established by \textcite{marcellino2006comparison}.

\textcite{plagborg2021local} were the first to actually show that under weak stationarity and unrestriced lag structure both methods estimate the identical impulse response function in population. Therefore any conventional LP impulse response can be obtrained through an appropriately ordered recursive VAR and any VAR impulse response through a LP with appropriate control variables. 
- they do not impose restrictions of the underlying DGP
The logic applies that any VAR model with sufficiently large lag lengths is able to capture the whole covariance properties of the data such that VAR($\infty$) forecasts coincide with direct LP forecasts. 
The authors come to the conclusion: 
- finite lag length, impulse response approximately agree if the same lag length is used for both but differ for $h \ge p$. 
- in population, linear local projections are exactly as robust to non-linearities in the DGP as linear VARs 

Here I will discuss the theoretical results that establish the equivalence of VARs and LPs under certain conditions, as shown by Plagborg-Moller \& Wolf (2021) and Fusari et al. (2026). I will also explain how this equivalence is affected by finite sample properties and dynamic misspecification, as discussed by Li et al. (2024), Montiel Olea et al. (2025), and Ludwig (2026).

The conventional finite-sample comparison between local projections and vector autoregressions is typically conducted at equal lag order, pitting LP(p) against VAR(p). Ludwig (2026) shows that this convention rests on a misleading notion of parsimony. His Theorem 1 establishes an exact finite-sample identity: the h-step prediction of an LP(p) is the forward iteration of a sequence of one-step VAR predictions whose lag order rises from p to $p + h - 1$. An equal-order VAR(p) is therefore a restricted special case of LP(p), obtained by holding the lag order fixed along the recursion rather than allowing it to increase. The familiar empirical pattern—lower approximation bias but higher sampling variance for LP relative to a same-order VAR—is, from this vantage point, largely a mechanical consequence of comparing a more general specification with its own restriction, rather than evidence of an intrinsic property of either estimator. This matters because it leaves the standard reading of the bias-variance trade-off, as articulated by Plagborg-Møller and Wolf (2021) and Montiel Olea et al. (2025), confounded: at equal lag order, the method effect (LP versus VAR) cannot be separated from the complexity effect (more versus fewer effective autoregressive parameters).
